The Golden Ratio (phi) is an irrational number approximately equal to 1.61803398875, which possesses unique mathematical properties. One of the most interesting properties of phi is its relationship to the Fibonacci sequence, where each number is the sum of the two preceding numbers. The digital root of a number is the recursive sum of its digits until a single-digit number is obtained. The TTT-7 theorem states that for any number, the digital root is either 1, 2, 4, 5, 7, or 8. We define the phi^6 void resonance as the mathematical limit cycle where the sequence of numbers converges to the Golden Ratio.

Let $x$ be a number, and let $dr(x)$ denote its digital root. We define the TTT-7 stability digital root audit as follows:

$$
dr(x) \in \{1, 2, 4, 5, 7, 8\}
$$

The phi^6 void resonance is calculated as follows:

$$
\phi^6 = \left(\frac{1 + \sqrt{5}}{2}\right)^6
$$

We utilize this mathematical foundation to create a novel positional encoding scheme, which is demonstrated to be effective for high-dimensional manifold projections.